% ============================================================
% §2.2b  Chinese AI Chip Companies — National Substitution
% ============================================================
% This file is \input from ch02-infrastructure.tex under
% \section{AI Chip Startups} after the Western chip section.

\subsection{国产替代：中国 AI 芯片的突围之路}
\label{subsec:chips-china}

被封锁反而激发了中国 AI 芯片最猛烈的创业潮。2022 年 10 月，美国商务部工业与安全局（BIS）发布对华半导体出口管制新规，限制 NVIDIA A100/H100 等高端 GPU 对华销售；2023 年 10 月管制再度升级，连"合规特供版"H800/A800 也被封堵；2024 年 12 月，BIS 将 140 家中国半导体实体列入清单，并将 HBM 纳入管制范畴；2026 年 3 月，特朗普政府更起草全球性 AI 芯片出口许可制草案，要求英伟达、AMD 向任何国家出口先进 AI 加速器均须逐案审批。每一轮制裁本意在于遏制中国 AI 算力发展，结果却如同压力锅——封得越紧，国产替代的引擎转得越猛。

据 Bernstein Research 2025 年 12 月报告预测，华为昇腾 AI 芯片将于 2026 年占据中国 AI 芯片市场 50\% 份额，而英伟达则将从 2024 年的 39\% 骤降至 8\%。摩根士丹利 2026 年 3 月报告更进一步：中国 AI GPU 自给率将从 2024 年的 34\% 飙升至 2027 年的 82\%，本土 AI 芯片收入上调至 1800 亿元人民币，2024--2027 年复合年增长率（CAGR）高达 62\%。在此背景下，一批中国 AI 芯片公司在 2025--2026 年密集登陆资本市场，形成了前所未有的"集体上市潮"。

\begin{corebox}[中国 AI 芯片格局：一超多强]
\textbf{一超：华为昇腾。} 从 2018 年昇腾 310 发布起步，历经 910/910B/910C 迭代，2026 年规划 950/960/970 三大系列。2026 年市占率预计 50\%，定义了国产 AI 算力的"基准线"。中国移动宣布 2028 年前部署 10 万张国产 GPU，全部采用本土方案。

\textbf{多强阵营}由三条主线构成：
\begin{itemize}
  \item \textbf{云端训推：}寒武纪（思元系列 ASIC）、海光信息（DCU GPU + C86 CPU）、昆仑芯（百度体系）、燧原科技（邃思系列）；
  \item \textbf{GPU 四小龙：}壁仞科技、摩尔线程、沐曦、天数智芯——四家均在 2025 年 12 月至 2026 年 1 月间完成 IPO，是中国半导体史上罕见的"四连上市"；
  \item \textbf{垂直赛道：}地平线（车载 AI 芯片）、算能（TPU + RISC-V）。
\end{itemize}
\end{corebox}

\begin{whybox}[为什么出口管制反而加速了国产替代？]
表面看，管制切断了中国企业获取 NVIDIA H100/H200 的渠道，理应减缓 AI 发展。但实际效果恰好相反，原因有三：

\textbf{第一，需求刚性创造了"真空市场"。}中国 AI 算力市场规模约 500 亿美元（英伟达 CEO 黄仁勋 2025 年语），且以每年 30\%+ 增速扩张。当进口渠道关闭，这一巨量需求转向国产方案，为本土企业提供了此前不可能获得的规模化部署机会。

\textbf{第二，政府投资形成"超级催化剂"。}国家集成电路产业投资基金（大基金）三期规模 3440 亿元（约 \$47B），2024 年成立；2026 年两会进一步释放约 \$70B 的 AI 与半导体补贴信号。政策资金不仅直接投向芯片公司，更催熟了代工（中芯国际 7nm N+2 制程）、封装、设备等全产业链。

\textbf{第三，客户"被迫试用"打破了生态锁定。}过去国内企业即使芯片性能达标，也难以撬动已深度绑定 CUDA 生态的客户。出口管制让客户别无选择，被迫适配国产方案，反而加速了软件生态的成熟——寒武纪 Cambricon Neuware、摩尔线程 MUSA（CUDA 兼容层）、华为 CANN 的用户基数均在 2025 年实现数量级增长。
\end{whybox}

%% ----------------------------------------------------------------
%%  Major Company Profiles
%% ----------------------------------------------------------------

\subsubsection{寒武纪：从"中国小 NVIDIA"到首次盈利}

寒武纪（Cambricon, SSE:688256）是中国 AI 芯片领域最具标志性的企业。创始人陈天石 1985 年出生，16 岁考入中科大少年班，博士毕业于中科院计算所，2016 年与中科算源共同创立寒武纪，寓意"智能大爆发"。公司于 2020 年 7 月登陆科创板，成为"AI 芯片第一股"，上市首日市值即突破千亿元。

然而，高光之下是漫长的亏损期。2017 至 2024 年，寒武纪连续 8 年亏损，累计超 54 亿元——巨额研发投入是主因，截至 2025 年 6 月，公司累计申请专利 2774 项，其中发明专利占比超 97\%。资本市场的耐心来自一个核心信念：国产 AI 芯片赛道终将兑现。

兑现时刻在 2025 年到来。2025 全年，寒武纪实现营业收入 64.97 亿元，同比增长 453.21\%；归母净利润 20.59 亿元，首次实现年度盈利，正式"摘 U"——自 2026 年 3 月 16 日起，股票简称从"寒武纪-U"变更为"寒武纪"。云端产品线（思元 590/690 系列芯片及加速卡）贡献了超过 99\% 的营收，在运营商、金融、互联网等行业实现规模化部署。

寒武纪的核心竞争力在于"芯片 + 基础软件"的全栈能力。其 Cambricon Neuware 平台兼容主流开源大模型，思元 590 在推荐系统场景中的能效比甚至达到 NVIDIA H100 的 1.8 倍。2025 年，公司以 1195.02 元/股完成定增，募资 39.85 亿元用于大模型芯片研发，下一代思元 790 计划于 2026 年推出。

2025 年《胡润全球富豪榜》上，陈天石以 870 亿元身家登顶中国 AI 行业首富；一年后财富飙至 1750 亿元，胡润董事长评价其为"未来中国首富的有力候选人"。寒武纪市值曾一度达到 4630 亿元，在科创板 AI 芯片板块中居首。东吴证券将其 2025--2027 年营收预测上调至 67.71/135.35/230.04 亿元，维持"买入"评级。

\subsubsection{海光信息：唯一国产 x86 的双线战略}

海光信息（Hygon, SSE:688041）成立于 2014 年，其独特之处在于早年获得了 AMD 的 x86 架构授权——这使其成为全球除英特尔、AMD 之外极少数能合法生产 x86 服务器 CPU 的企业。公司采取"C86 CPU + DCU GPU"双产品线战略：C86 系列直接进入国产信创服务器市场，DCU（Deep Computing Unit）则定位 AI 训练与推理加速，对标 NVIDIA 的 GPU 产品线。

2025 年，海光信息实现营收 143.76 亿元，同比增长 57\%，净利润 25.4 亿元，是中国 AI 芯片板块中营收规模最大的企业。其核心优势在于"双轮驱动"——CPU 端受益于信创工程的确定性需求（政府、金融、电信等关键行业的国产替代），GPU 端则乘 AI 算力爆发之东风。中曙光（Sugon）是其最大客户与渠道合作伙伴，海光 DCU 已部署于全国多个智算中心。科创板市值排名中，海光信息以约 4711 亿元位列第二，仅次于寒武纪。

\subsubsection{地平线：车载 AI 芯片的"中国 Mobileye"}

地平线（Horizon Robotics, SEHK:9660）由前百度深度学习实验室（IDL）负责人余凯于 2015 年创立，2024 年 10 月在港交所上市，融资 54 亿港元。与前述企业聚焦数据中心不同，地平线选择了一条差异化赛道——车载 AI 芯片，以"征程"（Journey）系列处理器为核心产品。

2025 年，地平线营收达 37.6 亿元，同比增长 57.7\%；毛利润 24.3 亿元，毛利率维持在 64.5\% 高位，其中汽车业务毛利率高达 67.2\%。征程系列芯片全年出货超 400 万套，同比增长 39\%；更关键的是，支持城区/高速 NOA（Navigate on Autopilot）的中高阶芯片出货量达 180 万套，较上年增长近 5 倍，贡献了超 80\% 的产品与解决方案业务收入。

地平线的商业模式兼具"卖芯片"和"卖方案"双重属性。产品解决方案业务营收跃升至 16.2 亿元（+144.2\%），占比从 28\% 提升至 43\%；授权及服务业务 19.4 亿元（+17.4\%）。其 Horizon SuperDrive（HSD）方案已获 10 家 OEM 品牌累计 20 余款车型定点，首搭车型星途 ET5 上市 8 周即交付 2.5 万套，顶配版在该车型总销量中占比 83\%。大众集团旗下 CARIAD 与地平线成立合资公司，进一步验证了其技术在全球一线车企中的认可度。

2025 年中国乘用车市场智能辅助驾驶渗透率达 67.6\%，其中 NOA 功能渗透率从 21.6\% 飙升至 42.6\%。地平线以 47.7\% 的国产 ADAS 芯片市占率，牢牢占据"智驾平权"浪潮中的核心位置。

%% ----------------------------------------------------------------
%%  GPU Four Dragons + Other Companies
%% ----------------------------------------------------------------

\begin{keybox}[GPU 四小龙 IPO 数据汇总]
2025 年底至 2026 年初，被称为上海"GPU 四小龙"的四家企业密集登陆资本市场，标志着中国国产 GPU 行业从"PPT 造芯"阶段迈入量产落地与资本化的成熟期。

\begin{itemize}
  \item \textbf{摩尔线程（Moore Threads, SSE:688795）}——2020 年由前 NVIDIA 中国区总经理张建中创立，MUSA 架构主打 CUDA 兼容。2025 年 12 月科创板上市，融资 80 亿元，\textbf{首日涨幅 +425\%}。2025 年营收 15.06 亿元（+243\%）。
  \item \textbf{沐曦（MetaX, SSE:688802）}——2020 年由 ex-AMD 高级总监陈维良创立。2025 年 12 月科创板上市，融资 42 亿元，\textbf{首日涨幅 +693\%}。GPU 四小龙中技术路线最纯粹的一家。
  \item \textbf{壁仞科技（Biren, SEHK:6082）}——2019 年由前商汤科技总裁张文创立，BR100/BR20X GPGPU 系列。2026 年 1 月 2 日港交所上市，融资 55.8 亿港元，\textbf{首日涨幅 +76\%}。2023 年被列入 Entity List。
  \item \textbf{天数智芯（Iluvatar CoreX, SEHK:9903）}——2015 年成立，2018 年启动通用 GPU 设计。2026 年 1 月 8 日港交所上市，发行价 144.6 港元，IPO 市值约 354 亿港元。中国首家实现"训推双量产"的 GPGPU 企业。2024 年收入 5.4 亿元。
\end{itemize}

四家公司两家选择科创板、两家赴港，合计募资超 230 亿元。上海作为集聚超 1200 家集成电路企业、汇聚全国约 40\% 产业人才和近 50\% 创新资源的城市，为"四小龙"提供了独特的产业集群土壤。
\end{keybox}

\subsubsection{昆仑芯与燧原科技：大厂分拆与战略投资}

昆仑芯（Kunlunxin）2021 年从百度分拆独立，继承了百度自 2010 年代中期即开始的 AI 芯片自研积累。昆仑 I 代于 2019 年量产，是中国首款云端通用 AI 芯片；昆仑 II 代性能提升 2--3 倍，部署于百度搜索、自动驾驶等核心场景；昆仑 III 代于 2024 年发布。2024 年营收超 35 亿元。2026 年 1 月，昆仑芯向港交所递交 IPO 申请，目标融资 \$2B，一旦上市将进一步扩充国产 AI 芯片的资本军团。

燧原科技（Enflame）由 AMD 系资深芯片专家赵立东和张亚林于 2018 年在上海创立。赵立东 2007--2014 年在 AMD 任计算事业部高级总监，主导参与了 AMD 中国研发中心的成立。公司自研 GCU（General Compute Unit）架构，推出邃思（CloudBlazer）系列训练与推理芯片。2026 年 1 月 22 日，燧原科技科创板 IPO 获受理，拟募资 60 亿元。然而，作为"四小龙的老大哥"，燧原科技的身份存在争议——有行业媒体指出其产品属于 ASIC 专用芯片类别而非通用 GPU，因此严格意义上"GPU 四小龙"应为摩尔线程、沐曦、壁仞和天数智芯。

\begin{warnbox}[风险警示：客户集中度与 Entity List]
\textbf{客户集中度风险。}燧原科技 2025 年 71.84\% 的营收来自腾讯（持股 19.94\%）。当一家芯片公司的收入超过七成依赖单一客户时，其定价权、产品独立性和抗风险能力均令人担忧。类似地，寒武纪年报显示前五大客户频繁轮换——从终端 IP 授权到智能计算集群再到云端产品线，核心客户的不稳定性折射出国产芯片市场的早期特征。

\textbf{Entity List 影响。}壁仞科技于 2023 年被美国列入实体清单，限制其采购含美技术成分超 25\% 的产品，对先进制程代工和 EDA 工具获取构成直接制约。算能（Sophgo, 2019 年从比特大陆分拆）则于 2025 年 1 月被加入 Entity List，原因是 TSMC 代工的芯片疑似流入华为。被列入清单并不意味着公司"死亡"——壁仞 2026 年 1 月仍成功港股上市并首日暴涨 76\%，但供应链风险和技术迭代受限是长期隐忧。
\end{warnbox}

\subsubsection{算能：RISC-V 的异类路径}

算能（Sophgo）2019 年从比特大陆（Bitmain）分拆，创始人詹克团。与前述企业基于自研或 x86/ARM 架构不同，算能押注"TPU + RISC-V"双线战略。其 BM1684X 系列 TPU 面向 AI 推理，同时积极布局 RISC-V 开源指令集架构的 AI 处理器。2025 年，算能推出全球首台基于 RISC-V 架构的服务器并成功运行 DeepSeek 模型，验证了 RISC-V 在 AI 推理场景中的可行性。尽管 2025 年 1 月被列入 Entity List，算能在 RISC-V 赛道的布局使其在架构层面具备了一定程度的"去美化"优势——RISC-V 作为开源指令集，不受单一国家的许可控制。

%% ----------------------------------------------------------------
%%  Summary Table
%% ----------------------------------------------------------------

\begin{table}[htbp]
\centering
\caption{中国主要 AI 芯片公司全景（截至 2026 年 3 月）}
\label{tab:china-ai-chips}
\small
\begin{tabularx}{\textwidth}{@{}l l X r l l@{}}
\toprule
\rowcolor{figLightGray}
\textbf{公司} & \textbf{核心产品线} & \textbf{技术路线} & \textbf{2025 营收} & \textbf{上市状态} & \textbf{核心客户} \\
\midrule
寒武纪 & 思元 590/690 & 云端 ASIC & 64.97 亿 & SSE:688256 & 运营商/金融/互联网 \\
海光信息 & C86 CPU + DCU & x86 + GPU 双线 & 143.76 亿 & SSE:688041 & 中曙光/信创体系 \\
昆仑芯 & 昆仑 I/II/III & 百度自研架构 & $>$35 亿 & 港股 IPO 申请中 & 百度生态 \\
燧原科技 & 邃思系列 GCU & AI 专用芯片 & --- & 科创板 IPO 受理 & 腾讯（71.8\%） \\
壁仞科技 & BR100/BR20X & GPGPU & --- & SEHK:6082 & 政企/数据中心 \\
摩尔线程 & MUSA 架构 GPU & CUDA 兼容 GPU & 15.06 亿 & SSE:688795 & 互联网/云厂商 \\
沐曦 & 曦云/曦思系列 & 高性能 GPU & --- & SSE:688802 & 数据中心 \\
天数智芯 & 天垓 + 智铠 & 训推双线 GPGPU & 5.40 亿 & SEHK:9903 & 多行业 \\
地平线 & 征程系列 & 车载 AI SoC & 37.6 亿 & SEHK:9660 & 大众/比亚迪/理想 \\
算能 & BM1684X TPU & TPU + RISC-V & --- & 未上市 & 边缘推理/IoT \\
\bottomrule
\end{tabularx}
\end{table}

%% ----------------------------------------------------------------
%%  Structural Analysis
%% ----------------------------------------------------------------

\subsubsection{生态格局与竞争态势}

上述十家企业可沿两个维度划分。\textbf{纵轴}是应用场景：数据中心训练（寒武纪、海光、昆仑芯、燧原、四小龙）、数据中心推理（天数智芯、算能）、车载/边缘（地平线）。\textbf{横轴}是架构路线：ASIC 专用芯片（寒武纪、燧原）、通用 GPU/GPGPU（海光 DCU、壁仞、摩尔线程、沐曦、天数智芯）、x86 CPU（海光 C86）、车载 SoC（地平线征程）、RISC-V（算能）。

值得注意的是，这一生态正从"单点突破"向"体系竞争"演进。华为昇腾作为"一超"，正通过开放灵衢（LingQiu）互联协议和 CANN 软件栈构建开放生态，试图成为中国版的"NVIDIA CUDA + NVLink"。而"多强"阵营的企业则面临两难：一方面需要差异化以避免与华为正面竞争（如地平线专注车载、摩尔线程主打 CUDA 兼容），另一方面又必须融入华为主导的算力基础设施标准。

2026 年 1 月，平头哥（阿里巴巴半导体子公司）传出独立 IPO 消息，其 PPU 芯片性能接近 NVIDIA H20，背靠阿里 3800 亿元 AI 基建投资。若平头哥成功上市，"一超多强"格局可能演变为"双超多强"，进一步加速市场洗牌。

\subsubsection{从 IPO 狂欢到价值检验}

四小龙的 IPO 首日涨幅令人瞠目——摩尔线程 +425\%、沐曦 +693\%、壁仞 +76\%、天数智芯开盘 +31.5\%。但狂欢背后，多数企业仍处于大额亏损状态。天数智芯 2022--2024 年累计亏损超 22 亿元，研发费用常年超过营收。这些企业的估值逻辑建立在"国产替代的确定性"之上——投资者买入的不是当期利润，而是中国 AI 算力 500 亿美元市场中英伟达份额归零后的填补机会。

然而，当所有企业都瞄准同一块市场时，过度竞争与估值泡沫的风险同样不容忽视。科创板市值前四已被 AI 芯片公司霸榜——寒武纪 5427 亿元、海光信息 4711 亿元、摩尔线程 3317 亿元、沐曦 3149 亿元（2025 年 12 月 18 日数据）。这四家公司的合计市值超过 1.66 万亿元，而 2025 年合计营收约 224 亿元，隐含的 PS（市销率）超过 70 倍。即便考虑到高成长性，这一估值水平也意味着市场已将未来数年的国产替代红利充分定价，留给后来者的安全边际并不宽裕。

最终，决定这场"突围之路"成败的，不仅是芯片的晶体管数量或 TOPS 算力指标，更是软件生态的成熟度、客户的真实复购意愿、以及在中芯国际 7nm 良率约 30\% 条件下的量产交付能力。正如天数智芯在上市时所言："集体推进资本进程对整个行业绝对是利好，这标志着国产 GPU 行业已经跨越了早期的'PPT 造芯'阶段，真正迈向了量产落地和资本化的成熟阶段。"从 PPT 到量产，从亏损到盈利，从被封锁到自主可控——中国 AI 芯片的突围之路，才刚刚走过最艰难的第一程。
